\chapter[Justice above the Law?]{Is There a Higher Justice above the Law?}
\section{A Two-Tablet Fridge Magnet}
Begin with a small object often sold in museum gift shops: a fridge magnet.

It resembles two grey stone tablets, rounded at the top, engraved with ten short lines imitating writing on stone over three thousand years ago. You find it in Jerusalem tourist shops, Washington museums, even bookshop gift counters. Not all buyers are religious; many simply recognise the image or like its appearance on a fridge.

Notice its message. Two tablets, rules handed down from above: do not kill, steal, bear false witness. Its core is not merely that rules matter, but something more specific and startling: \textbf{rules are not made by the powerful, but delivered to them}. The tablets rise above everyone's head, including the king's. Even a king stands beneath them.

Why remarkable? Through most human history, common sense said the opposite: law was the king's word. He originated, interpreted, and enforced rules. Telling him he had broken the law resembled telling a river it flowed wrongly. In many ancient societies, the sentence scarcely had a grammar.

Yet one tradition inscribed a king's capacity to offend and face judgement into its holiest books, reading them for over two millennia. This was the Hebrew tradition, ancient Israel's faith and writings, later central to Judaism and deeply entering Christianity and Islam.

The tablets hold our question: \textbf{is there higher justice above law?} If so, what is it, and who may declare it? When human rules conflict with it, which should an ordinary person obey?

First return to scenes where the question was sharply posed.

\section{A Lamb in the Royal Palace}
Come to Jerusalem around three thousand years ago, to King David's palace.

Imagine not a turreted European fairy-tale castle but thick-walled stone buildings high in a hill city, with small windows and courtyard daylight. At night, yellow olive-oil lamps mix smoke with flatbread and sheep's milk. Spear-bearing guards, document-keeping scribes, and people awaiting judgement fill the courtyard. In the ancient Near East, the king was the supreme court; judging was daily work.

An old man arrives: Nathan, a prophet. In Hebrew tradition, this means chiefly a spokesman for God, not a fortune-teller. Prophets address contemporary affairs, especially the powerful. Passing guards, he reaches David and offers neither accusation nor protest, only a story, broadly:

Two men lived in a town. One owned many flocks and herds; the other only a little ewe lamb raised by hand. It ate his bread, drank from his cup, slept in his arms, like a child. When a visitor came, the rich man spared his own animals and seized the poor man's sole lamb to slaughter for the guest.

According to Second Samuel in the Hebrew Bible, David was furious. With a ruler's judicial instinct, he declared the offender worthy of death and liable for fourfold repayment. He was acting as judge, indignant and decisive.

Then Nathan delivered one of Western literature's most famous accusations:

\begin{quote}
You are that man. (2 Samuel 12:7, Chinese Union Version)
\end{quote}
The background: David desired Bathsheba, wife of his officer Uriah, and committed adultery with her. To conceal it, he ordered Uriah into the most dangerous fighting, killing him through enemy weapons, then brought Bathsheba into the palace. Secret, legal, fully documented: troop orders followed procedure, battle deaths were normal, taking in a widow could even honour a fallen servant. Palace records showed no broken rule.

Nathan's lamb makes David's own judgement turn back upon him. We shall later examine the narrative's skill. For now remember the moment: the kingdom's most powerful man pinned by an unarmed elder's small story to his own verdict.

The account gives David one response: ``I have sinned against the Lord.'' No killing Nathan, excuses, or procedural imprisonment. Yet notice that he neither abdicated nor was deposed nor faced a human court. He remained king until death. Keep this crucial detail for later.

\section{Who May Judge the King?}
State the real difficulty before answering.

David's case knots itself: the king originates law, appoints judges, commands soldiers, receives final appeals. \textbf{When the lawmaker breaks the law, to whom can you appeal?}

Only a few options lie on the table, each troublesome.

One: law is the king's word, so royal acts are automatically lawful. Nathan's accusation then fails. Taking a subject's wife resembles a shepherd using his flock; no offence exists. Many ancient Near Eastern royal ideologies followed this route: the king represented deity on earth, his will was order. Its cost is toothless justice, governing ordinary people but not the person most needing restraint.

Two: higher law stands above enacted law and king alike. Nathan assumes this, citing no palace regulation but rules against adultery and murder, supposedly universally known though absent from palace files. Immediately: where is this higher law written? Who reads it? What if they misread? If several claim access and disagree, whom do we hear?

Three: no higher law exists. Nathan was merely brave and eloquent; David happened to listen. Another king and prophet produce another ending: a severed head and a palace entry punishing subversive falsehoods. Such endings may be more common. Is justice above law an illusion created by selecting surviving stories afterwards?

This is no palace antique. Its modern forms appear daily: rule-makers breaking school rules, oversight of platform moderation, obedience to a state's unjust law, refusal in conscience.

Before continuing, answer: Nathan commands no troops; David commands a kingdom. \textbf{How could a lamb story win?}

\section{Voices of Palace and Wilderness}
Hear at least two positions fully.

\textbf{First: the case for kingship.}

Strengthen the king's side. Repeated stories of criticised rulers make monarchy seem absurd from the outset, everyone awaiting a prophet's rebuke. Reality differs.

First Samuel says Israel initially lacked kings. Tribal leaders called judges addressed problems as they arose. Elders later demanded a king from the ageing prophet Samuel, explicitly wanting what neighbouring nations had: someone to govern, lead, and fight for them. Philistine military pressure weighed on the borders. The tribal alliance's improvised mobilisation repeatedly struggled.

Samuel disliked the request. He warned, broadly, that their king would requisition sons for chariots and runners, daughters for cooking and perfumes; take the best fields, vineyards, and olive groves for servants; claim a tenth of crops and livestock; and reduce them to servants. Later cries would come too late.

This warning in First Samuel 8 is unusual political literature: a regime introduces itself by reproducing its own indictment. More striking, the people hear it and still demand a king. Not stupidity, but the real necessity of safety and order. Judges repeatedly describes the kingless age:

\begin{quote}
In those days Israel had no king; everyone did as they pleased. (Judges 21:25, Chinese Union Version)
\end{quote}
That sounds free, but means no arbiter, no protection, the strong devouring the weak. Kingship offers unified judgement, organised defence, predictable order. Safety or freedom, order or risk: people chose three thousand years ago and still choose today.

\textbf{Second: the wilderness voice.}

Advance over two centuries to the eighth century BCE. Israel has divided into northern and southern kingdoms. Amos appears, not a religious professional, introducing himself:

\begin{quote}
I was neither a prophet nor a prophet's disciple; I was a herdsman and a tender of sycamore trees. (Amos 7:14, Chinese Union Version)
\end{quote}
A southern rural herder and fruit grower travels to northern Bethel, a religious centre, and denounces its busy markets and grand sanctuary. The wealthy ``sell the righteous for silver, and the poor for a pair of shoes'' (Amos 2:6, Chinese Union Version): creditors enslave the poor over tiny debts, perhaps a shoe's value, while bribed courts look away. He condemns merchants' undersized weights, ivory couches, bowls of wine, and indifference to Joseph's suffering.

The priest Amaziah warns him that this is the royal sanctuary and palace: speak at home, not here. Amos does not leave.

Notice his structure. He does not begin by accusing his own country. Amos opens with judgements on \textbf{foreign nations}: Damascus, Gaza, Tyre, Edom, Ammon, Moab, Judah, their atrocities counted before Israel's. This sequence carries an explosive point: his standard governs not only Israelites but everyone. It brings us to \textbf{universal ethics versus rules of a particular community}.

Community customs govern insiders: addressing elders, festival foods, Sunday rest. Across a village boundary they may cease to apply without difficulty. Amos's rules against murder and selling impoverished debtors into slavery supposedly govern all peoples, even nations disbelieving Israel's God. An early literary declaration: some justice recognises no border.

This higher law has a specific Hebrew-tradition name: \textbf{covenant}. God and people enter a reciprocal agreement, land and protection for observance of law. Alongside religious rites come social provisions: sabbatical debt release, field corners for poor gleaners, no interest charged to one's brothers. Crucially, \textbf{the king is also a contracting party, and the one most capable of breach}. Deuteronomy 17 requires future kings to copy the law and read it throughout life, lest their hearts rise above their brothers. Subjects are brothers, not property. Legitimacy derives not from bloodline or victories but obedience to the contract the king also signed.

The tradition thus makes a daring institutional design: kingship is necessary to prevent everyone doing as they please, yet kingship faces judgement beneath covenant. Together, these claims form its legacy.

\section[Thinking Bridge: A Rule's Higher Authority]{Thinking Bridge: Find the Authority above a Rule}
Leaving palace and wilderness, take a tool. Ask three questions of family, class, school, platform, or national rules:

\textbf{First: Who made the rule? (Source)}

Everyone together, or a few for the many? Different origins owe different explanations. A family's agreed lights-out time differs from a mother's unilateral announcement, even at the same hour.

\textbf{Second: Does it govern its maker? (Coverage)}

David's direct legacy. If teachers prohibit classroom phones, may theirs ring? If kings prohibit seizure, does taking Uriah's wife count? Exempt the maker and law quietly becomes a management tool. Tools act on others; law binds everyone, lawmakers included. Similar names, different substance.

\textbf{Third: If the rule is wrong, where can you appeal? (Appeal mechanism)}

Discuss class rules with teachers, school rules with heads; laws have amendment procedures, courts, public debate. A system answering every complaint with ``Rules are rules'' announces no higher court above it. That itself answers whether justice stands above law, negatively.

Together, these form a rule check-up. Add one conceptual distinction and engrave it in memory:

\textbf{Legality and justice are separate axes.}

Legality asks whether an action follows current rules; justice whether it is right. Their crossing creates four cells:

\noindent
\begin{tabularx}{\linewidth}{llX}
\toprule
 & Just & Unjust \\
\midrule
\textbf{Legal} & \shortstack[l]{Normal: pay taxes and\\ honour debts} & Lawful wrong: trading enslaved people under the laws then in force \\
\textbf{Illegal} & \shortstack[l]{Unlawful good: shelter\\ persecuted innocents\\ despite penalties} & Wrong on both counts: murder and robbery \\
\bottomrule
\end{tabularx}

\smallskip
Disputes often put one speaker on each axis. ``He broke a school rule'' and ``The rule was unfair'' can be \textbf{simultaneously true}, yet people argue for half an hour before discovering different coordinates. Ask whether the issue is legality or justice, and half the disputes dissolve.

\section{Two Prophetic Statements}
Read primary material to see higher justice's shape in this tradition.

Speaking for God, Amos rejects festivals, solemn assemblies, offerings, and songs, then declares:

\begin{quote}
Let fairness roll like mighty waters, and righteousness flow like a surging river. (Amos 5:24, Chinese Union Version)
\end{quote}
Its edge lies in context. It rejects not another religion but \textbf{its own people's procedurally correct observances}: calendared festivals, prescribed offerings, professional singing, all legal. Yet where a shoe-sized debt sells a poor person, such rites are worthless and repellent. Moral rightness judges religious correctness. This was a jarring ancient voice, when religions generally cared more for properly performed rites than fair market weights.

Another statement quoted for over two millennia comes from Micah:

\begin{quote}
Humanity, the Lord has shown you what is good. What does he ask of you? To act justly, love mercy, and walk humbly with your God. (Micah 6:8, Chinese Union Version)
\end{quote}
Notice the opening: humanity, not Israelites. The addressed requirement again crosses communal boundaries.

Compare something about a millennium earlier. Hammurabi carved his code on a black stele over two metres high, now in the Louvre. Its preface broadly says the gods appointed him to promote justice, destroy wrongdoing, and prevent strong from oppressing weak. Babylonian kings also wrote beautifully of protecting weakness. The crucial difference: \textbf{Hammurabi dispenses justice}; \textbf{prophetic kings undergo its scrutiny}. The vocabulary remains, the subject reverses. One king awards himself justice's medal; another is summoned before its court. Which requires more courage to preserve needs no explanation.

\section{Why Did Prophets Dare Rebuke Kings?}
With evidence, compare explanations. What happened: Israelite writings around the eighth century BCE preserve extensive prophetic criticism of rulers and wealth, a textual fact. Why this tradition produced it is interpretive and contested. Three main accounts follow.

\textbf{Explanation One: Within covenantal faith.}

Israel's identity rests on Exodus memory: a people claiming enslaved beginnings, a God siding with slaves and leading them from Pharaoh. Such divine character determines the contract. Protection for poor people, strangers, orphans, and widows is no optional virtue but explicit provision. Prophets invent no standard; they prosecute breaches, placing signed terms before offenders. This explains their \textbf{source of authority}. Amos claims not to advise but to read a contract older and higher than the king. Its limit: complete force only within the faith. Outsiders still require reasons to believe God sides with slaves. Distinguish believers' conviction that God made covenant, the tradition's claim that covenant outranks kingship, and the documented fact that these writings criticise kings. The first two are faith claims, the third textual evidence.

\textbf{Explanation Two: Social and economic conflict.}

Shift from heaven to earth. Eighth-century Israel faces expanding Assyria. Tribute, royal projects, armies require money; trade enriches urban merchants. Financial pressure descends onto rural society. Bad harvests force borrowing, mounting debt costs land, farmers become debt slaves. Amos's shoes describe concrete realities: bribed courts, debt mountains, ruined smallholders. Prophecy is bankrupt countryside accusing urban palaces. This grounded account explains concentration in \textbf{this period} and repeated attacks on ivory beds, summer houses, wine bowls. Its limit: prophets supply no economic programme, land reform, or revised debt bill; they denounce and leave. Calling them ancient social reformers projects a modern occupation backwards.

\textbf{Explanation Three: How the texts were made.}

Ask how these books reached us. Around 587 BCE, Babylon captured Jerusalem, destroyed the temple, and deported many elites. The state fell. Much prophetic writing was edited, organised, and finalised around that catastrophe. It serves to \textbf{explain disaster}. Why defeat? Covenant-breaking, oppression, other gods: not divine impotence but breach's cost. Prophets' prior warnings partly reflect survivors' retrospective arrangement. This sober account has textual grounding, yet cannot explain why traumatised people preserve \textbf{their own indictment} as sacred scripture. Defeated peoples often blame bad luck or overwhelming enemies; few canonise deserved punishment. Editorial timing does not fully explain treasured self-accusation.

Each holds some truth. Now an easily misjudged counterexample. You may have formed a comforting plot: brave weakness challenges power, justice wins, everyone rejoices. \textbf{Delete that template.} David confesses but neither abdicates nor faces trial. His first child with Bathsheba dies, interpreted as punishment. What does the king pay in human accounts? Family tragedy and a bowed head. Uriah remains dead. The tradition honestly preserves an awkward ending: the prophet wins the argument, the king retains his throne. Justice does not win; it is \textbf{spoken and remembered}. Read a fairy tale of inevitable right defeating wrong and you lose the sharpest point. It demands not belief in victory but someone speaking justice to kings and recording it even when it loses.

\section[Further Challenge: Is Unjust Law Law?]{Further Challenge: Is an Unjust Law Still Law?}
Now the heaviest question, with a technical name and twenty-four centuries of history.

\textbf{If a rule is procedurally lawful but profoundly unjust in content, is it still law, and do we owe it obedience?}

Two broad camps have argued without end.

\textbf{First: Law above law (the natural-law tradition).}

An early literary expression is Sophocles's \textit{Antigone}, staged in Athens around 441 BCE. Thebes has emerged from civil war. Two brothers killed each other; Polyneices brought foreign forces against his city. New ruler Creon forbids his burial as a traitor, under death penalty. His sister Antigone buries him, is caught, and says, broadly, that she obeys unwritten eternal laws, neither today's nor yesterday's, always alive with unknown beginnings. No mortal ruler's command overrides them.

The play's brilliance is refusing to make Creon a clown. \textbf{State his full case}: civil-war blood is barely dry. Bury a foreign-inviting traitor like a defender, and what is civic loyalty worth? Impartial law demands no exceptions to the burial ban. Antigone asks precisely for \textbf{an exception grounded in kinship}: he is her brother. Creon's cruelty and refusal to reconsider are wrong, but putting civic survival above family affection is not madness. Good sister versus bad king flattens tragedy. Its structure is justified visions of justice colliding and killing almost everyone. Higher law's danger is that everyone claims to hear it.

Two thousand years later, medieval Thomas Aquinas systematises this in the \textit{Summa Theologiae}: human law derives from higher eternal and natural law; an unjust human law is more violence than genuine law. Countless later disputes over obedience summon this statement from the shelf.

Nor is this exclusively Western. Confucian tradition contains an equally powerful charge. In ``King Hui of Liang, Part Two'' in the \textit{Mencius}, King Xuan of Qi asks whether Tang's exile of Jie and King Wu's campaign against Zhou were permissible killings of rulers by subjects. Mencius replies:

\begin{quote}
One who violates benevolence is a robber; one who violates righteousness is a destroyer. Such a robber and destroyer is a mere fellow. I have heard of executing the fellow Zhou, never of murdering a ruler. (\textit{Mencius}, ``King Hui of Liang, Part Two'')
\end{quote}
Broadly, trampling benevolence and righteousness reduces the tyrant to an isolated ordinary fellow. Wu killed that man, not a ruler. The logic: \textbf{moral qualifications define kingship; thorough injustice forfeits the identity itself}. Killing becomes removal of tyranny, not regicide. Amid repeated emphasis on ruler--subject bonds, this preserves a theoretical exit. Loyalty belongs not to the occupant's chair but to what it represents. Betray that and the chair is empty.

\textbf{Second: Law is law; goodness is another question (legal positivism).}

Give this side full strength. Calling it unjust law's accomplice is deeply unfair.

Its central claim: whether something is law and whether it is good law require \textbf{separate answers}. However bad, a duly enacted law remains law. Amend it through procedure rather than declare it nonexistent.

Why separate? First, \textbf{predictability}. If every judge and citizen can invalidate rules by personal justice, a thousand laws result and nobody knows what conduct is lawful. Order collapses, justice follows. Second, \textbf{who judges}? Each person hears higher law differently. Antigone does; Creon believes civic survival supreme too. Wars, persecution, and massacres claiming higher justice, divine will, or historical mission rival those claiming enacted law. Inquisitorial burning claimed precisely a law above human law. Third, \textbf{procedure is part of justice}. One successful bypass teaches everyone to bypass, removing the weak's last protective fence.

Their position can therefore be noble. Calling a bad law law defends not its content but everyone's inescapable task of amendment. Acknowledge it as law and accept the duty to change it; deny it and seemingly remove the matter from human procedure.

Remember the camps through two formulas. Natural law: \textbf{extreme injustice does not deserve the name law}. Positivism: \textbf{unjust law remains law, which is why it must change}. Their destination is close: neither advocates bowing to injustice. They dispute \textbf{routes}, individual refusal versus procedural reform, and \textbf{risks}: self-appointed higher courts versus unjust law endlessly surviving failed procedures under legality's cloak.

The twentieth century supplied a bloody footnote. After World War II, Nuremberg defendants repeatedly claimed obedience to orders and conformity with contemporary German law. The tribunal effectively answered that certain boundaries exist: crossing them is criminal regardless of lawful paperwork. Higher law returned heavily to modern law's centre.

\section[Rules, Reporting, and Following Procedure]{Today: School Rules, Reporting, and Just Following Procedure}
The three-thousand-year-old problem works beside you.

\textbf{Scene One: class.} One pupil misbehaves and the entire group loses conduct marks. You object: why punish me when I did not speak? The teacher invokes rules and collective honour. Apply the tool. Who made it? The teacher. Does it govern its maker? The teacher escapes collective punishment. Where can you appeal? If only ``No talking back,'' the classroom declares no higher court. One says you broke a rule, legality; the other that the rule is unjust. Both can hold. Childhood grievances are differently weighted versions of David's and Antigone's problem.

\textbf{Scene Two: a phone.} A platform removes a video for violating community standards. It writes, interprets, and usually hears appeals on those standards. Company rules are not laws, but unconstrained makers and no higher appeal resemble ancient palaces structurally. One difference: you can leave a platform; ancient people could not leave their kingdom. Hence special concern over inescapable platforms, whose rules approach legal power while accepting only management-tool obligations.

\textbf{Scene Three: history class.} In 1963, minister Martin Luther King Jr was arrested for anti-segregation protest and wrote from Birmingham jail. Alabama segregation was \textbf{legal}, supported by state statutes and court rulings. Defending lawbreaking, he invoked Augustine and Aquinas: unjust law is not law, and moral responsibility can demand disobedience. Notice his restraints: act \textbf{openly}, not through sabotage; \textbf{nonviolently}, harming nobody; \textbf{accept punishment}, using imprisonment to awaken public awareness; target specific injustice, not all law. Later called civil disobedience, this differs from ignoring rules through those conditions. Secretly running a red light serves convenience; openly occupying a prohibited seat and awaiting arrest submits a public indictment of the law.

Honesty includes the reverse. On the same land, opponents of vaccines and mask mandates and attackers of the legislature also claim conscientious resistance. Holocaust bureaucrat Eichmann defended himself in Jerusalem as faithfully following orders and regulations; the court rejected this, and he was executed in 1962. Together these expose the difficulty: \textbf{saints and villains alike have claimed obedience to higher justice and obedience to superiors}. Conscience alone cannot distinguish King from a legislative-building mob. Evidence, argument, openness, responsibility for consequences, and universally applicable principles rather than personal convenience can help, still subject to others' and time's scrutiny. Natural law and positivism keep arguing because each has seen the other's demons.

\section{Your Turn: The Rule Is Wrong. What Next?}
Return to the fridge magnet.

Tablets raised above a king mark one of civilisation's rare successes in making supreme power admit something above it. We have also seen the consequences. Someone must read the writing and may err or lie. Higher-law speakers may be Amos or crusaders; procedural defenders may protect order or prolong injustice.

Resize the question for yourself.

A rule in your group seems plainly unjust: collective deductions, biased award criteria, or a family double standard permitting adults what children cannot do. Three paths lie open.

First, obey while seeking amendment through authorised channels: suggestions, conversations, proposals. Slow, possibly fruitless, but preserving procedure.

Second, violate it openly and nonviolently, prepared for consequences, to reveal its injustice. Fast and visible, but personally costly, requiring justification for placing your judgement above the rule.

Third, outward compliance, private evasion. Easiest and commonest, but changing nothing and teaching hypocrisy.

Which, and why? Reweigh both sides. Creon guarded order, Antigone kinship and higher law; few won. Samuel warned, yet people chose kingship after tasting unregulated life. David confessed and remained king: the argument won, power did not lose.

A harder final question: if \textbf{you} become rule-maker, class leader, group leader, parent, platform moderator, how will you design rules so that someone like Nathan three thousand years later can say ``You are that man'' and you can hear it?

No standard answer. Remember the three questions: who made it, does it govern its maker, where can error be appealed? Next time someone says rules are rules, recognise a position, not discussion's endpoint.
